\begin{python0}
from solutions import *; clear()
\end{python0}
\sectionthree{Leaving things out}

Refer to the diagram from the previous section.

Actually it turns out that you can leave out \texttt{[stmt1]}. Try this
\begin{console}
for (int i = 0; i < 5; i++)
{   
    std::cout << "i have " << i
              << " head(s)"
              << std::endl;
}
\end{console}
And now modify it as follows:
\begin{console}
int i = 0;
for (; i < 5; i++)
{   
    std::cout << "i have " << i
              << " head(s)"
              << std::endl;
}
\end{console}
You can also leave the third part out as well!!!
\begin{console}
int i = 0;
for (; i < 5;)
{   
    std::cout << "i have " << i
              << " head(s)"
              << std::endl;
    i++;
}
\end{console}
What if the second part is blank? Try this:
\begin{console}
int i = 0;
for (;;)
{   
    std::cout << "i have " << i
              << " head(s)"
              << std::endl;
    i++;
}
\end{console}
(Pretty gross, isn't it? ...)

This tells you that if the boolean expression is missing, C++ will
assume you want the boolean value to be \EMPHASIZE{\texttt{true}}. Remember
what I said before: you have to pay attention whenever C++ does
something for you automatically.

By the way, the above program is executing an \EMPHASIZE{infinite loop}
(because it doesn't stop). This format of the for-loop is sometimes used in writing games, although we'll see that there's another type of loop called the
\EMPHASIZE{\texttt{while-loop}} that is frequently used instead. Although in the above

\tab[3em]{\texttt{for (;;)}}\\
\tab[3em]{...}

it seems like you can't get out of the loop. But, in fact, \EMPHASIZE{you
can get out of the loop in another way}. We'll talk
about it later.

